\subsection{Exact degree-six Euler factors for the exponent-nineteen quotient}
\label{subsec:nineteen-euler}

Retain $C:y^{19}=x(1-x)^7$, its rational order-three automorphism
$c$, and $A=\operatorname{Jac}(C/\langle c\rangle)$ from
Proposition~\ref{prop:lm-nineteen-rankone}.  Put
$K=\mathbb Q(\zeta_{19})$, $H=\langle7\rangle$,
$K_0=K^H$, and let $\sigma$ be the automorphism induced by
$\zeta_{19}\mapsto\zeta_{19}^2$.  The cyclic group
$\operatorname{Gal}(K_0/\mathbb Q)$ has order six, and its
complex conjugation is $\sigma^3$.  Write $\lambda_0$ for the
unique prime over nineteen in $K_0$.

\begin{lemma}[An exact auxiliary class number]\label{lem:nineteen-class-one}
The field $K_0$ has class number one.
\end{lemma}
\begin{proof}
Put $t=\zeta_{19}+\zeta_{19}^7+\zeta_{19}^{11}$. Its minimal
polynomial is
\[
 f(T)=T^6+T^5+2T^4-8T^3-T^2+5T+7,
\]
and the conductor--discriminant formula gives
$\operatorname{disc}K_0=-19^5$. Its Minkowski bound is
\[
 \left(\frac4\pi\right)^3\frac{6!}{6^6}\sqrt{19^5}
 <\left(\frac{200}{157}\right)^3\frac5{324}\,19^2
       \frac{4359}{1000}<51.
\]
Here $\pi>157/50$ and $\sqrt{19}<4359/1000$. The splitting
orders in the cyclic quotient of $(\mathbb Z/19\mathbb Z)^\times$
show that the prime ideals of norm at most fifty are exactly the
six over seven, the six over eleven, and the unique prime over
nineteen. Define
\[
\begin{split}
 u_7&=-\frac{18t^5+30t^4+56t^3-81t^2-72t+42}{77},\\
 u_{11}&=1-u_7,\\
 u_{19}&=\frac{-9t^5-15t^4-28t^3+79t^2-41t-21}{77}.
\end{split}
\]
Reduction modulo $f$ computes monic integral characteristic
polynomials for these three elements and gives
$N(u_7)=7$, $N(u_{11})=11$, and $N(u_{19})=19$.
Thus they are integral and generate primes of the indicated norms.
Galois conjugation generates all primes above seven and eleven;
the prime above nineteen is unique. Every ideal class therefore has
a representative which is a product of principal prime ideals,
by the displayed Minkowski bound. This proves the assertion.
All reductions, norms, and the exhaustive list of small prime ideals
are verified by the exact arithmetic certificate accompanying the paper.
\end{proof}

\begin{proposition}[The degree-six Hecke factors]\label{prop:nineteen-euler}
The threefold $A$ has CM by $K_0$, with its CM algebra defined over
$K_0$. Its algebraic Hecke character has conductor $\lambda_0$ and
CM type $\{0,1,2\}$ in the cyclic order determined by $\sigma$,
up to conjugation. Its rational conductor is $19^6$ and its root
number is $-1$. For the order-nine Dirichlet character
$\chi(2)=\zeta_9$, one has
\[
 L(J,s)=L(A,s)L(A\otimes\chi,s)L(A\otimes\bar\chi,s),
\]
where all three factors have conductor $19^6$, gamma factor
$(2\pi)^{-3s}\Gamma(s)^3$, and local Euler factor one at nineteen.
The Euler polynomials at every other prime and the exact complex
root number are given by the formulas proved below.
\end{proposition}
\begin{proof}
\par\medskip\noindent\textbf{The CM character and its normalization.} 
The endomorphisms $eae$, for $a\in K$, identify the corner
$e(K\rtimes H)e$ with $K_0$.  For $a\in K_0$ the endomorphism
$ea$ is defined over $K_0$, since the rational projector $e$
commutes with $a$ and $H$ fixes $a$.  Thus all the CM algebra
of $A$ is defined over $K_0$.  Its rank-one CM character,
denoted $\theta$, has values in $K_0$ and restricts to the
Jacobi-sum character $j$ over $K$.  This follows either by
applying the CM Hecke-character construction to $A$, or by
restricting the explicit representation of $J$ and applying
the projector $e$.  In the cyclic order determined by
$\sigma$, its type is $\{0,1,2\}$ up to rotation or inverse.
The inverse type is a rotation of this same set, so this
choice has no effect on the rational Euler polynomial.

For a prime $\ell\ne19$, let $e_\ell$ be its order modulo
nineteen and let $f_\ell$ be its order in
$(\mathbb Z/19\mathbb Z)^\times/H$.  If $\mathfrak p$ is a
prime of $K_0$ over $\ell$ and $\mathfrak P$ a prime of $K$
above it, then
\[
 \beta=\theta(\mathfrak p)\in K_0,\qquad
 j(\mathfrak P)=\beta^{e_\ell/f_\ell}.
\]
The quotient $e_\ell/f_\ell$ equals one for $e_\ell=1,2$
and three for $e_\ell=3,6,9,18$.  A cube root in $K_0$ is
unique, since $\mu(K_0)=\{\pm1\}$.

The restriction identity also proves the conductor descent,
without assuming an unstated CM conductor formula.  At
nineteen the extension $K/K_0$ is tame of degree three,
so its two absolute Galois groups have the same wild
inertia.  \citet[Proposition~3.12(1), author p.~16]{ShuExpanded2021}
gives conductor one for $j$; in particular it is trivial
on wild inertia and nontrivial on inertia.  The same two
assertions therefore hold for $\theta$, whose conductor
over $K_0$ is exactly $\lambda_0$.  Away from nineteen it
is unramified, since $A$ is a factor of the Jacobian of
a curve with good reduction away from nineteen.

We use the Jacobi convention in \citet[\S2.1, author p.~5]{ShuExpanded2021}: $j(\mathfrak P)=-\sum\chi(x)\chi(1-x)^7$.
In this convention
\[
 j(\mathfrak P)\equiv1\pmod{\lambda_K^2},
 \qquad \lambda_K=(1-\zeta_{19}).
\]
For an elementary verification, expand each character value
to first order in $\zeta_{19}-1$.  Over the residue field of
cardinality $q\equiv1\pmod{19}$, the sums of the character
exponents of $x$ and of $1-x$, with $x\ne0,1$, are divisible
by nineteen.  The constant term is $-(q-2)\equiv1\pmod{19}$.
This proves the congruence.  Stickelberger's identity and
the absolute value $|j|=\sqrt q$ are given in \citet[equation~(2.1), author p.~5]{ShuExpanded2021},
on the same page.  These three properties determine $j$:
two candidates with the same ideal and all absolute values
have root-of-unity quotient, and no nonidentity root of
unity in $K$ is congruent to one modulo $\lambda_K^2$.

\par\medskip\noindent\textbf{Primes splitting in $K_0$.} 
Suppose $\ell\bmod19\in H=\{1,7,11\}$, so $f_\ell=1$.
Take an algebraic integer $\alpha\in K_0$ of norm $\ell$;
its principal ideal is one of the six primes above $\ell$.
Set
\[
 b=\alpha\sigma(\alpha)\sigma^2(\alpha).
\]
Then $b\bar b=N_{K_0/\mathbb Q}(\alpha)=\ell$.
Replacing $\alpha$ by another generator of the same prime
changes $b$ by $\pm1$: the corresponding product of a unit
has absolute value one in every embedding, and hence is
a root of unity.  Since $\ell^3\equiv1\pmod{19}$, its
residue satisfies $b\equiv\pm\ell^2\pmod{\lambda_0}$.
Choose the unique sign such that
\[
 \beta\equiv\ell^2\pmod{\lambda_0}.
\]
For $e_\ell=1$, this gives $\beta\equiv1\pmod{\lambda_0}$;
for $e_\ell=3$, it gives $\beta^3\equiv1\pmod{\lambda_0}$.
In $K$ the valuation of $\lambda_0$ is three, so either
congruence implies the required congruence modulo
$\lambda_K^2$.  To check the ideal multiplicities exactly,
the nine-element reflex type in $K$ is the inverse image
of a three-element type in $K_0$.  When $e_\ell=1$, each
prime of $K_0$ splits into three primes of $K$, and the
nine factors in \citet[equation~(2.1), author p.~5]{ShuExpanded2021} consist of all three
primes over each of the three reflex primes of $K_0$.
Thus $(j(\mathfrak P))=(\beta)\mathcal O_K$, up to a
cyclic conjugation of the type.  When $e_\ell=3$, each
prime of $K_0$ remains inert in $K$; the three members
of $H$ fix the prime and repeat each reflex prime three
times.  Thus $(j(\mathfrak P))=(\beta^3)\mathcal O_K$.
These are precisely the two claimed ideal identities.
The normalization just proved identifies the elements,
up to the harmless cyclic conjugation.  Consequently
\[
 P_{\ell,A}(T)=N_{K_0/\mathbb Q}(1-\beta T).
\]
This formula needs only a verified norm-$\ell$ generator,
not a conditional class-group computation.  Lemma~\ref{lem:nineteen-class-one} establishes the existence
of such generators for all these primes.

\par\medskip\noindent\textbf{Order-nine primes.} 
Suppose $e_\ell=9$, so $f_\ell=3$.  The decomposition group
in $K_0$ is $\langle\sigma^2\rangle$, and $\beta$ belongs
to its fixed field $k=\mathbb Q(\sqrt{-19})$.  For example,
this follows because $\beta^3=j(\mathfrak P)$ is fixed by
that group and the cube root in $K_0$ is unique.  The
type $\{0,1,2\}$ meets the two decomposition orbits with
multiplicities two and one.  Hence, if $\mathfrak l$ is
one of the primes of $k$ above $\ell$, then
\[
 (\beta)=\mathfrak l^2\bar{\mathfrak l}
          =\ell\mathfrak l
\]
after choosing one of the conjugate labels.  The imaginary
quadratic prime $\mathfrak l$ remains inert in $K_0/k$,
and its prime in $K_0$ remains inert in $K/K_0$.
More explicitly, the nine-element reflex type in $K$
meets the two decomposition orbits with multiplicities
six and three.  \citet[equation~(2.1), author p.~5]{ShuExpanded2021} therefore gives
$(j)=\mathfrak P^6\bar{\mathfrak P}^{3}$.
Since $j=\beta^3$, descending these exponents gives
$(\beta)=\mathfrak l^2\bar{\mathfrak l}$, as claimed.
The imaginary
quadratic field $k$ has class number one: its Minkowski
bound is $(2/\pi)\sqrt{19}<3$, and two is inert, so every
ideal class has a representative of norm one.
Choose $\alpha\in\mathcal O_k$ with norm $\ell$ and set
$\beta=\pm\ell\alpha$.  Select the sign by
\[
 \beta\equiv\ell^6\pmod{\lambda_0}.
\]
This is possible and unique because $\beta^2\equiv\ell^3$
and $\ell^9\equiv1\pmod{19}$.  Its cube is congruent to
one, and the same ideal/weight/congruence argument
identifies the correct Frobenius eigenvalue.  Thus
\[
 P_{\ell,A}(T)=1-\operatorname{Tr}_{k/\mathbb Q}(\beta)T^3
                   +\ell^3T^6.
\]
For example the formula gives $1+5T^3+125T^6$ at five
and $1+119T^3+4913T^6$ at seventeen.  Each can be computed
using a quadratic norm equation; no finite field of
cardinality $\ell^9$ is required.

\par\medskip\noindent\textbf{Even-order primes.} 
If $e_\ell$ is even, complex conjugation belongs to the
decomposition group.  The Jacobi sum is therefore real,
and its absolute values and normalization give
$j(\mathfrak P)=-\ell^{e_\ell/2}$.  Taking the unique
root when necessary gives $\beta=-\ell^{f_\ell/2}$ and
\[
 P_{\ell,A}(T)=
 \begin{cases}
 (1+\ell T^2)^3,&e_\ell=2\text{ or }6,\\
 1+\ell^3T^6,&e_\ell=18.
 \end{cases}
\]
These cases and the preceding two paragraphs exhaust all
primes other than nineteen.

\par\medskip\noindent\textbf{Three small completed functions.} 
Let $\chi$ be a primitive Dirichlet character of order nine
modulo nineteen.  It restricts nontrivially to $H$, whereas
a character of order three does not.  Induction in stages
and the extensions of $\theta|_{G_K}$ to $G_{K_0}$ give
\[
 V(J)\simeq V(A)\oplus V(A)\chi\oplus V(A)\chi^2.
\]
The representation $V(A)$ has the self-twist $\chi^3$,
since this character is trivial on $G_{K_0}$.  Thus the
last summand can also be written $V(A)\chi^{-1}$, and
\[
 L(J,s)=L(A,s)L(A\otimes\chi,s)L(A\otimes\bar\chi,s).
\]
The local twisted polynomials are simply
$P_{\ell,A}(\chi(\ell)T)$.  All three functions have
conductor $19^6$, gamma factor $(2\pi)^{-3s}\Gamma(s)^3$,
and local Euler factor one at nineteen.  To see the
conductor assertion, the Jacobi character over $K$ has
conductor exponent one at nineteen and a nontrivial
quadratic restriction to inertia there, while $K/\mathbb Q$
is tame.  Thus $V(J)$ is tame with no inertia invariants.
Its three dimension-six summands have no invariants
either, so each Artin conductor exponent is six.  The
order-nine twists are even and have the same infinity
type.  The root numbers of the two conjugate twists are
complex conjugates and have product one.  Since $J$ has
root number minus one, the root number of $A$ is minus one.
It is therefore sufficient to certify
\[
 L'(A,1)\ne0,\qquad L(A\otimes\chi,1)\ne0.
\]
Their product gives the desired nonzero derivative of
$L(J,s)$ at one.  The central nonvanishing and derivative are certified in
Proposition~\ref{prop:nineteen-analytic-certificate}.

\par\medskip\noindent\textbf{An explicit root number for the complex twist.} 
There is also an exact finite Gauss sum for the root number,
so it need not be estimated from a truncated functional
equation.  Normalize $\chi(2)=\zeta_9$, consistent with
the Euler twist $P_{\ell,A}(\chi(\ell)T)$.  With
$\zeta_{19}=\exp(2\pi i/19)$, the formula is
\[
 w(A\otimes\chi)=
 \frac{i}{\sqrt{19}}\sum_{u=1}^{18}
       \left(\frac{u}{19}\right)\chi(u)^6\zeta_{19}^{6u}.
\]
Here is a derivation including the local convention.
Use the positive finite additive character in \citet[\S3.1, equation~(3.4), author p.~9]{ShuExpanded2021} and its compatible
negative infinite character.  The complex local factors
are $i^{-1}$, as recorded in \citet[\S3.3, author p.~13]{ShuExpanded2021};
the three complex places therefore contribute $i$.
The trace additive character of $K_{0,\lambda_0}$ has
conductor $\lambda_0^{-5}$, since the degree-six local
extension is tame.  The multiplicative character has
conductor $\lambda_0$, so the integration shell in \citet[equation~(3.2), author p.~9]{ShuExpanded2021} is
$19^{-1}\mathcal O_{K_0,\lambda_0}^{\times}$.

Let $\varphi$ be the unitary character associated with
$\theta$.  On rational units $u\in\{1,\ldots,18\}$,
\[
 \varphi_{\lambda_0}(u)=\left(\frac{u}{19}\right).
\]
Indeed, for positive rational $u$ prime to nineteen,
restriction and ideal norm give
$\theta((u))^3=j((u))=u^9(u/19)$.
The last equality follows from the tame quadratic local
character in \citet[equation~(3.14), author p.~16]{ShuExpanded2021}, or
directly from the ideal, absolute value and Jacobi
congruence.  Taking the unique cube root in $K_0$ gives
$\theta((u))=u^3(u/19)$, and the principal-idele relation
gives the asserted local value.  The local component of the
Dirichlet twist is the inverse of its residue character,
and the residue norm from $K_0$ is $u\mapsto u^6$.
Thus the twisted local character on these units is
$(u/19)\chi(u)^{-6}$.  Its value at nineteen is one:
all other finite local factors on this principal idele
are one, and the infinite unitary factors at the positive
rational number nineteen are one.

On the integration shell, write $x=u/19$.  The residue
additive factor is $\zeta_{19}^{6u}$, since
$\operatorname{Tr}_{K_0/\mathbb Q}(u/19)=6u/19$.
The self-dual volume of the local integer ring is
$19^{-5/2}$.  Multiplication by $19^{-1}$ changes volume
by $19^6$, while $|x|^{-1/2}=19^{-3}$; dividing into
nineteen residue classes leaves the common factor
$19^{-1/2}$.  Substitution in equation~(3.2) therefore
gives precisely the displayed sum after multiplying by
the infinite factor $i$.
For the trivial twist, $(6/19)=1$ and the quadratic Gauss
sum is $i\sqrt{19}$, recovering $w(A)=-1$.  For the
inverse twist the root number is the complex conjugate,
as follows either from the functional equation or directly
from this formula using the oddness of the quadratic
character.  These checks also fix the sign convention.

\end{proof}

\begin{lemma}[The published canonical quartic]\label{lem:nineteen-canonical}
The genus-three quotient has the following plane-quartic
model, already published as $X_{15}$ by
\citet[author pp.~2 and 22; Table~1, case~15, p.~5]{KilicerEtAl2018}.
Its identification with this Fermat quotient is the known isomorphism of
\citet[Remark~7.2, p.~12]{Tafazolian2026}:
\[
\begin{split}
0={}&a^4-a^3b+2a^3d+2a^2bd+2a^2d^2-2ab^2d\\
 &+4abd^2-b^3d+3b^2d^2+2bd^3+d^4.
\end{split}
\]
\end{lemma}
\begin{proof}
The holomorphic differentials on $C$ are
\[
 \omega_h=(1-x)^{\lfloor7h/19\rfloor}\frac{dx}{y^h},\qquad
 h\in\{5,8,10,12,13,15,16,17,18\}.
\]
Their orders at the three branch points are respectively
$18-h$, $19b+18-7h$, and $8h-19b-20$, where $b=\lfloor7h/19\rfloor$;
all are nonnegative. Direct substitution gives
$c^*\omega_h=(-1)^{b-h}\omega_{7h\bmod19}$.
The three orbit sums form the invariant basis. Thus the invariant
holomorphic differentials have common
factor $dx/y^{18}$ and coefficient functions
\[
\begin{split}
 a&=(1-x)y^{13}+(1-x)^5y^2-(1-x)^6y,\\
 b&=(1-x)^2y^{10}+(1-x)^6+(1-x)^4y^6,\\
 d&=(1-x)^3y^8-(1-x)^4y^5+(1-x)^5y^3.
\end{split}
\]
Substitution and reduction by $y^{19}=x(1-x)^7$ prove
the quartic relation.  The three differentials are linearly independent, as their coefficient
functions have disjoint powers of $y$ of degree less than nineteen.
They form the full canonical system of the genus-three quotient.
Its image lies in the displayed smooth quartic, whose smoothness
is part of the cited published model. Since a plane quartic has degree
four and the canonical divisor of a genus-three curve has degree
four, the canonical map has degree one. This identifies the quotient
with the published curve. Exact substitution and rational linear
algebra also verify that its six quadratic canonical monomials are
independent. Small-prime plane counts at seven and eleven are eighteen
and eight, agreeing with traces $-10$ and $4$ from the CM formula.
No general assertion about reduction of this particular integral
quartic is required by the coefficient algorithm.
\end{proof}
